\section{Regression}

\subsection{Least-Squares Linear Regression}

\begin{defbox}
Fit $y=mx+b$ so that all data points lie as close as possible to the line (small residuals). Error $=$ real $-$ predicted $= y_i-(mx_i+b)$. Squaring the errors prevents $+/-$ cancellation:
\[
SSE = \sum_{i=1}^n \big(y_i - mx_i - b\big)^2.
\]
\end{defbox}

\begin{examplebox}
\textbf{Proof of the least-squares formulas.}

\emph{Intercept.} $\dfrac{\partial\,SSE}{\partial b} = -2\sum(y_i-mx_i-b) = 0 \Rightarrow \sum y_i - m\sum x_i - nb = 0$, so
\[
b = \frac{\sum y - m\sum x}{n}. \tag{A}
\]

\emph{Slope.} $\dfrac{\partial\,SSE}{\partial m} = -2\sum x_i(y_i-mx_i-b) = 0 \Rightarrow \sum x_iy_i - m\sum x_i^2 - b\sum x_i = 0$. Substitute $b$ from (A):
\[
\sum xy - m\sum x^2 - \frac{\sum y - m\sum x}{n}\sum x = 0.
\]
Multiply by $n$ and collect terms in $m$:
\[
n\sum xy - mn\sum x^2 - \sum y\sum x + m\left(\sum x\right)^2 = 0
\]
\[
n\sum xy - \sum x\sum y = m\left[n\sum x^2 - \left(\sum x\right)^2\right]
\]
\[
\boxed{m = \dfrac{n\sum xy - \sum x\sum y}{n\sum x^2 - (\sum x)^2}} \qquad\blacksquare
\]
\end{examplebox}

\begin{examplebox}
\textbf{Worked simulation — fit a line and find $r$.} Data $(1,2),(2,4),(3,5),(4,4),(5,5)$, $n=5$.

\begin{center}
\begin{tabular}{@{}cccccc@{}}
\toprule
 & $x$ & $y$ & $x^2$ & $y^2$ & $xy$ \\
\midrule
& 1 & 2 & 1 & 4 & 2\\
& 2 & 4 & 4 & 16 & 8\\
& 3 & 5 & 9 & 25 & 15\\
& 4 & 4 & 16 & 16 & 16\\
& 5 & 5 & 25 & 25 & 25\\
\midrule
$\Sigma$ & 15 & 20 & 55 & 86 & 66\\
\bottomrule
\end{tabular}
\end{center}
\[
m = \frac{5(66)-15(20)}{5(55)-15^2} = \frac{330-300}{275-225} = \frac{30}{50} = 0.6
\]
\[
b = \frac{20-0.6(15)}{5} = \frac{20-9}{5} = \frac{11}{5} = 2.2 \qquad\Rightarrow\qquad \hat y = 0.6x+2.2
\]
\[
r = \frac{5(66)-15(20)}{\sqrt{[5(55)-15^2][5(86)-20^2]}} = \frac{30}{\sqrt{50\times 30}} = \frac{30}{\sqrt{1500}} = \frac{30}{38.7298}=0.7746
\]
$r^2 = 0.60$ — the line explains 60\% of the variance in $y$.
\end{examplebox}

\begin{qabox}
\textbf{Q1:} What does the correlation coefficient $r$ measure, and how do you read its value? \\
\textbf{A:} The strength and direction of the linear relationship between $x$ and $y$. $r=\pm1$: perfect positive/negative line; $|r|$ near $0.6$–$0.8$: fairly strong but scattered; $r=0$: no linear relation. Closer to $\pm1$ = stronger.

\textbf{Q2:} What does $r^2$ mean? \\
\textbf{A:} The fraction of variance in the data explained by the model; $r^2=1$ means 100\% explained, $r^2=0.6$ means 60\%.

\textbf{Q3:} Does a strong correlation imply causation? \\
\textbf{A:} No. Two variables can move together by coincidence, or because a hidden third factor drives both.

\textbf{Q4:} Which running sums must be tabulated to compute $m,b,r$? \\
\textbf{A:} $\sum x,\ \sum y,\ \sum x^2,\ \sum y^2,\ \sum xy$.
\end{qabox}

\subsection{Multiple Linear Regression}

\begin{defbox}
Model with $n$ features — a hyperplane in $n$ dimensions:
\[
\hat y = w_0 + w_1x_1+w_2x_2+\dots+w_nx_n,
\]
where $w_0$ is the bias/intercept and $w_1,\dots,w_n$ are weights. (1D $\to$ point, 2D $\to$ line, 3D $\to$ plane, higher $\to$ hyperplane. Without $w_0$ the fit is forced through the origin.)
\end{defbox}

\begin{examplebox}
\textbf{Derivation of the MSE gradient.} With $\hat y_i = \sum_k w_k x_{ik}$ (and $x_{i0}=1$) and cost
\[
J(\mathbf w) = \frac{1}{2m}\sum_{i=1}^m (\hat y_i-y_i)^2,
\]
\[
\frac{\partial J}{\partial w_j} = \frac{1}{2m}\sum_{i=1}^m 2(\hat y_i-y_i)\frac{\partial \hat y_i}{\partial w_j} = \frac{1}{m}\sum_{i=1}^m (\hat y_i-y_i)\,x_{ij}.
\]
Hence the concrete update:
\[
w_j \leftarrow w_j - \frac{\alpha}{m}\sum_{i=1}^m (\hat y_i-y_i)\,x_{ij}. \qquad\blacksquare
\]
\end{examplebox}

\begin{qabox}
\textbf{Q1:} List the five assumptions of multiple linear regression. \\
\textbf{A:} Linearity; Independence (samples independent); Homoscedasticity (constant error variance); Normality (errors $\sim\mathcal N$); No multicollinearity (features not linearly dependent on each other).

\textbf{Q2:} Scenario: the errors follow a Poisson distribution. Can you still apply MLR? \\
\textbf{A:} No — Poisson errors are not normally distributed, so the Normality assumption is violated; MLR is not appropriate here.

\textbf{Q3:} What does a violated homoscedasticity assumption look like? \\
\textbf{A:} The error variance is not constant — the spread of residuals grows, or some points sit extremely far from the fit (fanning out), instead of scattering evenly.

\textbf{Q4:} What does ``no multicollinearity'' require, and why does it matter? \\
\textbf{A:} The input features must not be linearly dependent on one another (not redundant); violating it makes individual weight estimates unstable/uninterpretable.

\textbf{Q5:} Why does the cost function use $\frac{1}{2m}$ rather than just $\sum(\hat y_i-y_i)^2$? \\
\textbf{A:} Squaring stops $+/-$ residuals from cancelling; the $\tfrac{1}{2m}$ factor removes scale effects (the 2 cancels the derivative's factor of 2, and dividing by $m$ averages over samples) so GD updates stay well-scaled.
\end{qabox}
